% !PRINTED 409
\weberpage{409}
%<<CONT
\gesperrt{unabhängigen Variablen $x_0, x_1, \ldots, x_{n-1}$, so kann
man}
\weq{4}{u = f(\vartheta),\ u_1 = f_1(\vartheta), \ldots, u_{\nu-1} = f_{\nu-1}(\vartheta)}
\gesperrt{setzen, worin $f, f_1, \ldots, f_{\nu-1}$ rationale Functionen
von $\vartheta$ sind und}
\weq{5}{\vartheta = \chi(u, u_1, \ldots, u_{\nu-1}) = \psi(x_0, x_1, \ldots, x_{n-1})}
\gesperrt{eine rationale Function der $u$, oder auch der
$x$ ist.}

Nach dem Satze 1. nämlich können wir zunächst, wenn wir
$x_0, x_1, \ldots, x_{n-1}$ durch rationale Functionen einer Variablen $t$
ersetzen, so dass $z$ nicht von $t$ unabhängig wird, die Functionen
$u_h$ durch die Formeln (4), (5) darstellen. Die Relationen
\[
u_h = f_h(\vartheta) = f_h[\chi(u, u_1, \ldots, u_{\nu-1})]
\]
sind dann in Bezug auf $t$ identisch befriedigt, und müssen also
nach dem Satze 2. auch in den Variablen $x$ identisch sein,
w. z. b. w.

4.~\gesperrt{Wenn von zwei Variablen $\vartheta, \vartheta_1$ jede eine ratio\-
nale Function der anderen ist, so sind sie lineare
Functionen von einander.}

Zum Beweise nehmen wir an, es sei
\weq{6}{\vartheta = \frac{\varphi(\vartheta_1)}{\psi(\vartheta_1)},\ \vartheta_1 = \frac{\varphi_1(\vartheta)}{\psi_1(\vartheta)},}
und verstehen unter $\varphi, \psi$ zwei ganze Functionen ohne gemeinsamen
 Theiler, ebenso unter $\varphi_1, \psi_1$. Ordnen wir die zweite der
Gleichungen (6) in der Form $\varphi_1(\vartheta) - \vartheta_1\psi_1(\vartheta) = 0$ nach Potenzen
 von $\vartheta$, so mag sie die Gestalt annehmen:
\[
a_0\vartheta^m + a_1\vartheta^{m-1} + \cdots + a_{m-1}\vartheta + a_m = 0,
\]
worin die Coëfficienten $a_0, a_1, \ldots, a_m$ \gesperrt{ganze lineare Functionen}
 von $\vartheta_1$ sind. Substituiren wir darin für $\vartheta$ den Werth
aus der ersten Gleichung (6), so folgt die in Bezug auf $\vartheta_1$ identische
 Gleichung
\[
a_0\varphi^m + a_1\varphi^{m-1}\psi + \cdots + a_{m-1}\varphi\psi^{m-1} + a_m\psi^m = 0.
\]
Hiernach muss $a_0\varphi^m$ durch $\psi$ theilbar sein, und weil $\varphi$
und $\psi$ relativ prim vorausgesetzt sind, so muss $a_0$ durch $\psi$
theilbar, also $\psi$ constant oder linear sein. Ebenso schliessen
wir, dass $\varphi$ constant oder linear sein muss, und da nicht beide
%CONT>>
